% Informe TP IA 2026 - formato LNCS (Springer).
% Requiere llncs.cls (y splncs04.bst si se usa BibTeX): bajarlos de la plantilla oficial
% http://preview.springer.com/gp/computer-science/lncs/conference-proceedings-guidelines
% y dejarlos en esta carpeta. Limite del enunciado: 12 carillas.
\documentclass{llncs}

\usepackage[utf8]{inputenc}
\usepackage[T1]{fontenc}
\usepackage[spanish]{babel}
\usepackage{graphicx}

\begin{document}

\title{Comparaci\'on entre \'Arboles de Decisi\'on (C4.5) y Random Forest\\en problemas de clasificaci\'on}

% TODO: completar autores, instituci\'on y mails
\author{Integrante 1 \and Integrante 2 \and Integrante 3 \and Integrante 4}
\institute{Universidad Tecnol\'ogica Nacional -- Inteligencia Artificial, 5\textordmasculine{} Nivel ISI, 2026\\
\email{mail@ejemplo.com}}

\maketitle

\begin{abstract}
% TODO: 150-250 palabras: problema, m\'etodo, datos, resultado principal.
\keywords{\'Arbol de decisi\'on \and C4.5 \and Random Forest \and clasificaci\'on \and ensambles.}
\end{abstract}

\section{Introducci\'on}
% TODO: contexto, objetivo del trabajo, organizaci\'on del documento.

\section{Descripci\'on de los algoritmos}
\subsection{Decision Tree (C4.5)}
% TODO: entrop\'ia, ganancia, raz\'on de ganancia, atributos continuos, poda.
\subsection{Random Forest}
% TODO: bagging, selecci\'on aleatoria de atributos, error de generalizaci\'on.

\section{Descripci\'on de la soluci\'on implementada}
% TODO: arquitectura del notebook, clase del \'arbol, interfaz y problemas encontrados
% (el enunciado pide explicarlos expl\'icitamente).

\section{Simulaciones y resultados}
% TODO: dataset, metodolog\'ia de evaluaci\'on, tablas y figuras.

\section{Conclusiones}
% TODO

\begin{thebibliography}{1}
% TODO: verificar cada referencia contra la fuente original antes de entregar.
\bibitem{liu}
Liu, B.: Web Data Mining: Exploring Hyperlinks, Contents, and Usage Data. 2nd edn. Springer.

\bibitem{breiman}
Breiman, L.: Random Forests. Machine Learning 45, 5--32 (2001)
\end{thebibliography}

\end{document}
